\documentclass[11pt]{letter}
\usepackage{geometry}
\geometry{margin=1in}
\usepackage{setspace}
\usepackage{parskip}

\begin{document}

\begin{letter}{
Editor-in-chief: Huimin Zhao
ACS Synthetic Biology\\
American Chemical Society}

\opening{Dear Editor,}

Please find enclosed our manuscript entitled:

\textbf{``Reservoir computing with coupled genetic oscillators for the Classification of Arrhythmia''}

submitted for consideration for publication in \textit{ACS Synthetic Biology}.

\textbf{Corresponding Author:} James Newsome \\
\textbf{Co-Authors:} Gizem Buldum, Tim Rudge etc.

\textbf{Motivation}

We believe this research is well-suited to the ACS Synthetic Biology journal because it bridges synthetic biology and computational neuroscience by proposing a novel use of synthetic genetic oscillators as computational substrate within a reservoir computer for machine learning classification. This work introduces a biologically inspired architecture capable of processing physiological time-series data, showcasing a promising direction for hybrid biological-computational systems. Our approach not only simulates coupled genetic oscillators to display the inherent intelligence within biological systems, but also leverages these properties for practical biomedical applications inspired by para-medicine, such as arrhythmia classification using electrocardiograms scans, thus demonstrating how synthetic biology can intersect with machine learning and medical diagnostics.

\textbf{Lay Summary}

This study explores how synthetic genetic circuits, specifically biological oscillators that mimic gene activity, can be used as the core of a computational system to identify irregular heartbeats (arrhythmias). Inspired by how cells regulate gene expression over time, akin to circadian rhythms, we designed a simulated network of genetic oscillators that respond dynamically to time-series data from real patient electrocardiograms. These biological-inspired oscillations are processed through a structure called a reservoir computer, which helps identify patterns in complex signals. By applying this system to arrhythmia classification, we demonstrate a new way that synthetic biology can be used for machine learning, and subsequently support medical diagnostics. The work opens the door to using living or bio-inspired systems as computational substrate, connecting biology and computing in a unique and impactful way.

We appreciate your time and consideration of our submission. Please do not hesitate to contact us with any questions.

\closing{Sincerely,\\
\vspace{1em}
James Newsome \\
Newcastle University, School of Computing \\
james.newsome02@gmail.com}

\end{letter}

\end{document}
